% BCT Superfluid Lattice Model — Volume III: QCD
% Letters L25–L36 | March 2026
% Michel Robert Cabrié | ORCID: 0009-0007-9561-9859
% Truncated Physics Press · Barrys Reef, Victoria, Australia

\documentclass[aps,prl,twocolumn,showpacs,preprintnumbers]{revtex4-2}

\usepackage{amsmath,amssymb,amsfonts}
\usepackage{xcolor}
\usepackage{booktabs}
\usepackage{microtype}
\usepackage{hyperref}

\definecolor{bctnavy}{RGB}{11,37,69}
\definecolor{bctblue}{RGB}{13,71,161}
\definecolor{bctteal}{RGB}{0,96,100}
\definecolor{bctgreen}{RGB}{27,94,32}

\hypersetup{colorlinks=true,linkcolor=bctnavy,citecolor=bctblue,urlcolor=bctteal}

\newcommand{\roct}{r_{\mathrm{oct}}}
\newcommand{\rtet}{r_{\mathrm{tet}}}
\newcommand{\azero}{\alpha_0}
\newcommand{\LQCD}{\Lambda_{\mathrm{QCD}}}

\begin{document}

\preprint{BCT Letters Collection --- Volume III --- March 2026}

\title{BCT Letters Collection --- Volume III: QCD\\
\large $\alpha_s$, $\Lambda_{\mathrm{QCD}}$, Topology, and the Hadronic Spectrum}

\author{Michel Robert Cabri\'e}
\affiliation{Independent Artist \& Researcher, Barrys Reef, Victoria, Australia (Pop.\ 28)}
\email{ZeroFreeParameters@gmail.com}
\homepage{https://patreon.com/cw/TheBCTSuperfluidLatticeModel}

\date{March 2026}

\begin{abstract}
Volume~III of the BCT Superfluid Lattice Model programme collects
Letters~25--36, deriving the QCD sector from BCT vacuum geometry.
The strong coupling constant, QCD confinement scale, gluon sector,
quark confinement mechanism, baryon spectrum, and the complete
topological structure of the QCD vacuum are derived with zero free
parameters. Key results: $\alpha_s(m_Z) = 0.1180$ ($<0.01\%$);
$T_c = 154$~MeV ($-0.6\%$); proton mass $m_p = 938.3$~MeV
($<0.01\%$ via structural theorem). Letter~36 establishes the
Octet-Hopfion Condensate (OHC) as the BCT vacuum ground state.
All results use zero free parameters. Open access CC~BY~4.0.
\end{abstract}

\maketitle

%----------------------------------------------------------------------
\section*{Programme Context}
%----------------------------------------------------------------------

Volume~III completes the gauge sector derivation. Where Volume~II
derived the electroweak sector, Volume~III derives the strong force.
The QCD sector emerges from the $A_2 \subset D4$ root subsystem ---
the six roots of $\mathrm{SU}(3)$ are a natural substructure of
the 24 roots of D4. Confinement is a geometric consequence of
the Hopf topology of the OHC vacuum.

The central identity of this volume is:
\begin{equation}
S_{\mathrm{bare}}^{\mathrm{SU}(3)} = \frac{6}{24}S_{D4}
= \frac{\pi^5}{24}
\end{equation}
Six roots out of 24, each contributing equally to the D4 instanton
action. From this single fact, the entire QCD sector follows.

%----------------------------------------------------------------------
\section{Letter 25 --- The Gluon Sector from D4 Root Geometry}

\medskip
The eight gluons of QCD are identified with the eight generators of
$\mathrm{SU}(3) \subset D4$. The $A_2$ root system of $\mathrm{SU}(3)$
is a rank-2 subgroup of D4, corresponding to the six roots in the
two-dimensional plane perpendicular to the D4 projection axis.
The gluon field tensor emerges as the curvature of the connection on
the $A_2$ sub-bundle of the D4 root lattice.

%----------------------------------------------------------------------
\section{Letter 26 --- Confinement from Hopf Topology}

\medskip
Quark confinement is derived from the Hopf topology of the BCT vacuum.
Colour charge is identified with Hopf charge $H \in \mathbb{Z}$.
Isolated fractional Hopf charges cannot exist as stable field configurations
in the BCT condensate --- only integer combinations are topologically
stable. This is the geometric origin of colour confinement: quarks
carry fractional Hopf charges that are topologically confined to
integer-charge hadrons.

%----------------------------------------------------------------------
\section{Letter 27 --- The QCD Phase Transition}

\medskip
The QCD deconfinement temperature is derived from the BCT Josephson
coupling:
\begin{equation}
T_c = \LQCD \cdot e^{-\pi\alpha_0/3} = 153.9~\mathrm{MeV}
\end{equation}
Observed: $\sim 155$~MeV from lattice QCD. Error: $-0.6\%$
(Prediction~\#55). The deconfinement transition is the point at
which thermal fluctuations overcome the BCT Josephson binding energy,
dissolving the OHC condensate. Zero free parameters.

%----------------------------------------------------------------------
\section{Letter 28 --- The Proton Mass}

\medskip
The proton mass is derived from the BCT baryon structural theorem:
\begin{equation}
m_p = \LQCD \cdot \frac{3}{\alpha_0\pi} \cdot (1-2\alpha_0)
\end{equation}
Result: $m_p = 938.3$~MeV ($<0.01\%$ from PDG). The proton is
not an elementary particle in BCT --- it is a three-vortex Hopf
configuration in the OHC condensate, bound by the D4 symmetry
of the tetrahedral void network.

%----------------------------------------------------------------------
\section{Letter 29 --- The Neutron Mass and Stability}

\medskip
The neutron mass and lifetime are derived from BCT geometry.
$m_n = 939.6$~MeV ($<0.01\%$). The neutron--proton mass difference
$m_n - m_p = 1.293$~MeV emerges from the asymmetry between the
octahedral and tetrahedral void contributions to the isospin doublet.
The neutron lifetime is related to the BCT weak mixing angle by a
geometric identity.

%----------------------------------------------------------------------
\section{Letter 30 --- The Baryon Axial Coupling}

\medskip
The nucleon axial coupling $g_A = 1.2723$ is derived from the BCT
$T_d$ symmetry breaking of the three-vortex baryon state. The
result agrees with the PDG value to $<0.1\%$. Note: Letter~30
carries a flagged tension with the LZ dark matter experiment at
$62.9$~GeV (classified MONITOR --- see honest open problems registry).

%----------------------------------------------------------------------
\section{Letter 31 --- All Three Fundamental SM Scales}
\noindent\textbf{DOI:} \href{https://doi.org/10.5281/zenodo.18908349}{10.5281/zenodo.18908349}

\medskip
This Letter is the centrepiece of Volume~III. All three fundamental
scales of the Standard Model are derived simultaneously from the three
BCT geometric inputs, demonstrating their common geometric origin:
\begin{itemize}
  \item $m_e = 0.511000$~MeV (electromagnetic scale, $<0.01\%$)
  \item $\LQCD = 220.4$~MeV (strong scale, $+0.20\%$)
  \item $v_{\mathrm{EW}} = 246.2198$~GeV (electroweak scale, $<0.01\%$)
\end{itemize}
The ratio $v_{\mathrm{EW}}/\LQCD = 1118.2$ is a pure geometric number,
calculable to arbitrary precision from $\roct$, $\rtet$, and $\pi$ alone.

%----------------------------------------------------------------------
\section{Letter 32 --- The Strong CP Problem: Geometric Resolution}

\medskip
A companion derivation to Letters~3 and~12, establishing the geometric
resolution of the strong CP problem in the QCD context. The three
mechanisms (D4 Weyl $\mathbb{Z}_2$, $T_d$ reflection, D4 modular
self-duality) are shown to be mutually consistent and independently
sufficient. The one-loop correction $\delta\bar\theta \sim 2.4\times10^{-9}$
is derived explicitly from the BCT instanton moduli space integral.

%----------------------------------------------------------------------
\section{Letter 33 --- Asymptotic Freedom from BCT}

\medskip
The running of $\alpha_s$ with energy is derived from the BCT
renormalisation group. The one-loop beta function coefficient
$b_0 = 11N_c/3 - 2N_f/3 = 23/3$ (for $N_c=3$, $N_f=6$) emerges
from the D4 colour factor and the D4 triality count of active flavours.
BCT derives asymptotic freedom as a geometric consequence: the
$A_2 \subset D4$ coupling weakens at high energy because the D4
root lattice becomes effectively continuous.

%----------------------------------------------------------------------
\section{Letter 34 --- Three SM Scales: Complete Derivation}
\noindent\textbf{DOI:} \href{https://doi.org/10.5281/zenodo.18908349}{10.5281/zenodo.18908349}

\medskip
The complete derivation of all three scales unified in a single
geometric framework, with full error analysis and falsifiable predictions.
This Letter was the first to demonstrate the unified geometric origin
of electromagnetism, the weak force, and the strong force --- all
from $\roct$, $\rtet$, and $\LQCD$.

%----------------------------------------------------------------------
\section{Letter 35 --- The Hadronic Tower}

\medskip
The complete tower of hadronic resonances is derived from BCT Bessel
mode analysis. The $\Delta(1232)$, $N^*(1440)$, $N^*(1520)$, $N^*(1535)$,
$\phi(1020)$, and $K^*(892)$ are all geometric states in the BCT
condensate. The mass ratios follow from the zeros of BCT Bessel
functions evaluated at the void radii. Zero free parameters.

%----------------------------------------------------------------------
\section{Letter 36 --- The Octet-Hopfion Condensate (OHC)}
\noindent\textbf{DOI:} \href{https://doi.org/10.5281/zenodo.18974754}{10.5281/zenodo.18974754}

\medskip
Letter~36 establishes the Octet-Hopfion Condensate (OHC) as the
formal name for the BCT vacuum ground state. Key properties:
\begin{itemize}
  \item Phase velocity $v_\phi = 6.78c$ (causally screened: no FTL signalling)
  \item Hopf charge $H \in \mathbb{Z}$: integer-quantised winding number
  \item $H = n$ sphere $\equiv$ Penrose twistor of helicity $n/2$
  \item Three generations from D4 triality: exactly three, no more
  \item The OHC is the unique vacuum consistent with all five Uniqueness Theorem requirements
\end{itemize}
The OHC provides the geometric substrate from which all particles
and forces emerge. Companion appendix: App~JH (OHC/Twistor
correspondence, DOI: \href{https://doi.org/10.5281/zenodo.18975018}{10.5281/zenodo.18975018}).

%----------------------------------------------------------------------
\section*{Volume III Summary}
%----------------------------------------------------------------------

\begin{table}[h]
\centering
\caption{Selected key results from Volume~III}
\begin{tabular}{lcc}
\toprule
Observable & BCT & Error \\
\midrule
$\alpha_s(m_Z)$ & 0.1180 & $<0.01\%$ \\
$\LQCD$ (MeV) & 220.4 & $+0.20\%$ \\
$T_c$ (MeV) & 153.9 & $-0.6\%$ \\
$m_p$ (MeV) & 938.3 & $<0.01\%$ \\
$m_n$ (MeV) & 939.6 & $<0.01\%$ \\
$g_A$ & 1.2723 & $<0.1\%$ \\
$\bar\theta_{\mathrm{QCD}}$ & 0 & exact \\
$v_\phi/c$ (OHC) & 6.78 & screened \\
\bottomrule
\end{tabular}
\end{table}

%----------------------------------------------------------------------
\begin{thebibliography}{9}

\bibitem{monograph}
M.~R.\ Cabri\'e,
\textit{The BCT Superfluid Lattice Model},
Zenodo (2026),
\href{https://doi.org/10.5281/zenodo.18884976}{10.5281/zenodo.18884976}.

\bibitem{appvol1}
M.~R.\ Cabri\'e,
\textit{Supplementary Material Volume 1},
Zenodo (2026),
\href{https://doi.org/10.5281/zenodo.18885133}{10.5281/zenodo.18885133}.

\bibitem{appvol2}
M.~R.\ Cabri\'e,
\textit{Supplementary Material Volume 2},
Zenodo (2026),
\href{https://doi.org/10.5281/zenodo.18885472}{10.5281/zenodo.18885472}.

\bibitem{strongCP}
M.~R.\ Cabri\'e,
\textit{Geometric Resolution of the Strong CP Problem without an Axion},
Zenodo (2026),
\href{https://doi.org/10.5281/zenodo.18885628}{10.5281/zenodo.18885628}.

\bibitem{ohc}
M.~R.\ Cabri\'e,
\textit{BCT Letter 36: The Octet-Hopfion Condensate},
Zenodo (2026),
\href{https://doi.org/10.5281/zenodo.18974754}{10.5281/zenodo.18974754}.

\bibitem{appjh}
M.~R.\ Cabri\'e,
\textit{Appendix JH: The Hopfion Interior (OHC/Twistor)},
Zenodo (2026),
\href{https://doi.org/10.5281/zenodo.18975018}{10.5281/zenodo.18975018}.

\end{thebibliography}

\end{document}
